\documentclass{article}

% if you need to pass options to natbib, use, e.g.:
%     \PassOptionsToPackage{numbers, compress}{natbib}
% before loading neurips_2026

% The authors should use one of these tracks.
% Before accepting by the NeurIPS conference, select one of the options below.
% 0. "default" for submission
% \usepackage{neurips_2026}
% the "default" option is equal to the "main" option, which is used for the Main Track with double-blind reviewing.
% 1. "main" option is used for the Main Track
%  \usepackage[main]{neurips_2026}
% 2. "position" option is used for the Position Paper Track
%  \usepackage[position]{neurips_2026}
% 3. "eandd" option is used for the Evaluations & Datasets Track
 % \usepackage[eandd]{neurips_2026}
% 4. "creativeai" option is used for the Creative AI Track
%  \usepackage[creativeai]{neurips_2026}
% 5. "sglblindworkshop" option is used for the Workshop with single-blind reviewing
 % \usepackage[sglblindworkshop]{neurips_2026}
% 6. "dblblindworkshop" option is used for the Workshop with double-blind reviewing
 \usepackage[dblblindworkshop]{neurips_2026}
% 7. "education" option is used for the Education Track
%  \usepackage[education]{neurips_2026}


% After being accepted, the authors should add "final" behind the track to compile a camera-ready version.
% 1. Main Track
 % \usepackage[main, final]{neurips_2026}
% 2. Position Paper Track
%  \usepackage[position, final]{neurips_2026}
% 3. Evaluations & Datasets Track
 % \usepackage[eandd, final]{neurips_2026}
% 4. Creative AI Track
%  \usepackage[creativeai, final]{neurips_2026}
% 5. Workshop with single-blind reviewing
%  \usepackage[sglblindworkshop, final]{neurips_2026}
% 6. Workshop with double-blind reviewing
%  \usepackage[dblblindworkshop, final]{neurips_2026}
% Note. For the workshop paper template, both \title{} and \workshoptitle{} are required, with the former indicating the paper title shown in the title and the latter indicating the workshop title displayed in the footnote.
% For workshops (5., 6.), the authors should add the name of the workshop, "\workshoptitle" command is used to set the workshop title.
\workshoptitle{Who Verifies the Agents? Toward Reliable Agent Development}
% 7. Education Track
% \usepackage[education, final]{neurips_2026}

% "preprint" option is used for arXiv or other preprint submissions
 % \usepackage[preprint]{neurips_2026}

% % to avoid loading the natbib package, add option nonatbib:
% \usepackage[nonatbib]{neurips_2026}

\usepackage[utf8]{inputenc} % allow utf-8 input
\usepackage[T1]{fontenc}    % use 8-bit T1 fonts
\usepackage{hyperref}       % hyperlinks
\usepackage{url}            % simple URL typesetting
\usepackage{booktabs}       % professional-quality tables
\usepackage{amsfonts}       % blackboard math symbols
\usepackage{amsmath}
\usepackage{nicefrac}       % compact symbols for 1/2, etc.
\usepackage{microtype}      % microtypography
\usepackage{xcolor}         % colors
\usepackage{graphicx}

% Note. For the workshop paper template, both \title{} and \workshoptitle{} are required, with the former indicating the paper title shown in the title and the latter indicating the workshop title displayed in the footnote. 
\graphicspath{{../plots/}{plots/}{./}}

\title{Agreement Is Not Verification:\\
Why LLM Judges Need the Environment, Not Each Other}


% The \author macro works with any number of authors. There are two commands
% used to separate the names and addresses of multiple authors: \And and \AND.
%
% Using \And between authors leaves it to LaTeX to determine where to break the
% lines. Using \AND forces a line break at that point. So, if LaTeX puts 3 of 4
% authors names on the first line, and the last on the second line, try using
% \AND instead of \And before the third author name.


\author{
  Arushi Waddepalli \\
  IIIT Lucknow \\
  \texttt{lit2024009@iiitl.ac.in} \\
  \And
  Johnny Kozman \\
  University of Hertfordshire \\
  \texttt{jg25abs@herts.ac.uk} \\
  \And
  Tilak Parajuli \\
  Tribhuvan University \\
  \texttt{tilak.765421@bkmc.tu.edu.np} \\
}

\begin{document}


\maketitle


\begin{abstract}
  %The abstract paragraph should be indented \nicefrac{1}{2}~inch (3~picas) on both the left- and right-hand margins. Use 10~point type, with a vertical spacing (leading) of 11~points. The word \textbf{Abstract} must be centered, bold, and in point size 12. Two line spaces precede the abstract. The abstract must be limited to one paragraph.%
  Agent pipelines increasingly close the loop with an LLM verifier: one model judges
  whether another's output is correct, and that judgment gates deployment, training, or
  search. We report a controlled study of that judgment across 64{,}800 verification
  calls, crossing four models as both generators and verifiers with three authorship
  frames and three verification strategies over three task domains. Two fixes commonly
  proposed for unreliable LLM judges both fail. The first is to blind the judge to
  authorship. Self-preference is real, but it is not about the self: telling a model it
  wrote an answer changes nothing, while a model that genuinely wrote it approves its own
  errors 12 percentage points more often. Nor is that gap self-recognition---rewriting an
  answer into another model's prose leaves the verdict unmoved, and in 9 of 12
  model$\times$domain cells what looks like recognition is a model agreeing with its own
  earlier conclusion rather than identifying its own writing. Self-preference is belief
  persistence, so blinding cannot remove it. The second fix is a better or larger panel.
  Errors from stronger generators are systematically harder to catch---the same direction
  in all three domains---and no verification strategy narrows the gap; where the task
  admits no executable check, judges fail together rather than independently, so adding
  judges buys nothing. Only an environment signal changes the picture, and even then
  judges contradict that signal on one call in eight, with differential testing supporting
  fewer than one in ten of those contradictions. The binding constraint on agent
  verification is the signal the environment can supply, not the model reading it.
\end{abstract}

\section{Introduction}
\label{Introduction}

Language models are increasingly asked to check other models' work. An LLM verifier
decides whether a candidate answer is correct, and that decision is load-bearing: it
filters training data, scores rollouts under AI feedback, gates a step in an agent's plan,
and closes the loop in self-refinement
pipelines~\citep{zheng2023judging,bai2022constitutional,madaan2023selfrefine}. The
verifier stands in for a correctness oracle that is too slow or simply unavailable, and
whether that substitution is safe depends on a property verifiers are rarely tested for:
that the verdict tracks correctness, and nothing else.

It does not. A verifier is itself a prompted model, sensitive to who it is told wrote an
answer, to how the answer is phrased, and to whether the answer agrees with what it would
have said~\citep{zheng2023judging,wang2023fair,panickssery2024llm}. Practice has converged
on two fixes. The first treats self-preference as a conflict of interest: blind the judge
to authorship, or never let a model grade its own output. The second treats unreliability
as noise: use a stronger judge, a better prompt, or a panel of
judges~\citep{verga2024replacing}. Neither has been tested under a design that separates
what a judge \emph{believes} about authorship from who actually wrote the text, or that
holds a specific error fixed while varying the judge that sees it.

We run that design. Four instruction-tuned models---Qwen2.5-72B, DeepSeek-V3,
Llama-3.3-70B and Mistral-Nemo-12B---each answer every item, and all four then verify
every candidate under three authorship frames (told it is their own, told it is another
model's, told nothing) crossed with three strategies (verdict only, brief chain of
thought, rubric). Crossing the told frame with the true generator--verifier identity
separates the belief of authorship from the fact of it; having every judge see every
candidate lets us hold a specific wrong answer fixed and ask which judges catch it,
removing item difficulty as a confound by construction rather than by adjustment. The
three domains vary one property deliberately---whether the environment can supply an
independent check. Code admits execution against tests, mathematics admits answer
matching, graduate-level science admits neither. In total, 64{,}800 verification calls
over 450 items.

Self-preference survives the intervention aimed at it, because that intervention targets
the wrong mechanism. Telling a model it wrote an answer moves false approval by under a
percentage point, in the direction opposite to self-preference, while a model that
\emph{actually} wrote the answer approves its own errors twelve points more often, holding
the error fixed. Nor is that gap self-recognition: rewriting an answer into another
model's prose, against a placebo rewrite, leaves the verdict unmoved, and in almost every
model and domain we test, apparent recognition dissolves once we control for the model
agreeing with the conclusion it originally reached. This runs against
\citet{panickssery2024llm}, where self-recognition drives self-preference. The mechanism we
find is belief persistence---a model that still holds a conclusion approves any answer
expressing it---so blinding a judge to authorship cannot remove a bias that does not run
through authorship.

The remaining unreliability is structural rather than correctable by better judging.
Errors from stronger generators are systematically harder to catch, in the same direction
in all three domains, and no strategy narrows the gap; where no mechanical check exists,
judges fail together rather than independently, so a panel inherits a shared blind spot
instead of averaging it away. What changes the picture is a signal from the environment
rather than a better model reading the answer: judge accuracy runs at 87.5\% where the
judge is shown an execution result, against 57\% to 65\% where it must judge unaided, and
the correlation between judges' errors disappears in the grounded case. Even then the
judge is not the source of that reliability---it contradicts the runtime result on one
call in eight, and independent differential testing supports fewer than one in ten of
those contradictions.

\paragraph{Contributions.}
(i) A protocol separating believed from actual authorship, with an error-stratified
estimator that holds the specific error fixed across judges. (ii) Three probes---style
transfer, structural preference, and self-recognition with a consistency control---that
jointly rule out self-recognition as the mechanism behind self-preference. (iii) Evidence
that error detectability falls with generator strength and that judge failures are
correlated, which together explain why stronger judges, better prompts and larger panels
all fail on ungrounded tasks. (iv) Evidence that the binding constraint on agent
verification is the signal the environment supplies, not the judge reading it. We release
all prompts, code, and the 64{,}800 verification traces.

\paragraph{Related work.}
LLM-as-a-judge evaluation is now standard practice, alongside a catalogue of its failure
modes: position and verbosity bias~\citep{zheng2023judging,wang2023fair}, and a preference
for a model's own generations that has been attributed to
self-recognition~\citep{panickssery2024llm}. Proposed remedies include stronger or more
numerous judges~\citep{verga2024replacing} and richer prompting. A parallel line finds that
models cannot reliably correct their own reasoning without external
feedback~\citep{huang2024selfcorrect}, and that execution- or process-level signals
substantially improve verification~\citep{liu2023evalplus,lightman2024verify}. We differ in
three ways: we separate believed from actual authorship rather than confounding them; we
test the self-recognition account against a consistency control; and we vary the
availability of an environment signal while holding the judge pool fixed, which lets us
attribute reliability to grounding rather than to judge capability.

\section{Setup and judging system}
\label{sec:setup}

\paragraph{Tasks.} We draw 150 items from each of three benchmarks with a fixed sampling
seed: HumanEval+~\citep{liu2023evalplus} for code, MATH~\citep{hendrycks2021math} for
mathematics, and GPQA-Diamond~\citep{rein2024gpqa} for graduate-level science. The domains
are chosen not for difficulty but for what the environment can verify, and Table~\ref{tab:domains}
records the distinction the paper turns on. Correctness is established by running the
candidate in code, by matching a final answer in mathematics, and by a multiple-choice
letter in science---but only in code does any of that reach the judge. Mathematics is
therefore a checkable task presented to an unaided judge, which separates
\emph{grounding} from \emph{checkability}. Science options are shuffled per item under a
fixed seed and rendered as A--D; items whose options fail a validity check are discarded
before sampling.

\begin{table}[t]
  \caption{The three domains differ in whether a mechanical check exists and, separately,
  in whether the judge is shown one. Only code verifications receive an environment signal.}
  \label{tab:domains}
  \centering
  \small
  \begin{tabular}{llll}
    \toprule
    Domain & Benchmark & Correctness established by & Signal shown to judge \\
    \midrule
    Code    & HumanEval+     & Execution against the test harness      & Execution result block \\
    Math    & MATH           & Final-answer extraction, numeric match  & None \\
    Science & GPQA-Diamond   & Option letter; ambiguous parses fail    & None \\
    \bottomrule
  \end{tabular}
\end{table}

\paragraph{Models.} Four instruction-tuned models act as both generators and verifiers:
Qwen2.5-72B-Instruct, DeepSeek-V3, Llama-3.3-70B-Instruct-Turbo, and
Mistral-Nemo-Instruct-2407. All are served through a single inference provider at
temperature~0. The pool deliberately spans a wide capability range, from a 12B dense model
to a 671B-parameter mixture of experts; \S\ref{sec:results-detect} exploits that spread,
and \S\ref{sec:limitations} notes it as a confound for any model-level ranking.

\paragraph{The verification grid.} Every model produces a candidate for every item, with
no generation failures (1{,}800 candidates), and every model then verifies every candidate.
Each verification varies two factors independently. The \emph{authorship frame} tells the
judge the answer is its own (\textsc{self}), another model's (\textsc{other}), or says
nothing (\textsc{neutral}); prompts are identical apart from that single sentence, and no
model is ever named. The \emph{verification strategy} requests a bare verdict
(\textsc{direct}), a brief chain of thought before the verdict (\textsc{cot}), or a
rubric-scored assessment (\textsc{rubric}). This yields $4\times4\times3\times3 = 144$
verifications per item and 64{,}800 in total.

The design's essential property is that the told frame is crossed with the true
generator--verifier identity. A judge told ``you wrote this'' is usually being misinformed,
and a judge told nothing is sometimes grading its own work. Comparing frames at fixed true
authorship isolates the \emph{belief} of authorship; comparing true authorship at a fixed
frame isolates the \emph{fact} of it. Section~\ref{sec:results-belief} reports both.

\paragraph{Execution grounding in code.} Code verifications, and only code verifications,
carry a block reporting what happened when the candidate was actually run against the
benchmark's harness: exit status, assertion count, captured output, and the final traceback
frame on failure. The prompt states that this is a signal rather than the final word and
explicitly permits the judge to mark passing code incorrect on logical grounds. That
permission is deliberate: it makes judge-initiated deviation from the environment signal
observable, which \S\ref{sec:results-ground} measures. The harness, its timeouts, and the
differential fuzzer used to adjudicate disputed overrides are described in
Appendix~\ref{app:execution}.

\paragraph{Verdicts and metrics.} Verifiers return JSON and a verdict is read from a single
boolean field, with fallbacks for fenced or prefixed output. Of 64{,}800 calls, 277
(0.43\%) never yielded a parseable verdict and are dropped; all of them come from the two
reasoning strategies, and 262 from the smallest model. Approving an incorrect answer is a
\emph{false approval} and rejecting a correct one a \emph{false rejection}. We report raw
accuracy for comparability with prior work but treat balanced accuracy as primary: the rate
at which a judge catches errors is maximised by rejecting everything, and raw accuracy is
inflated wherever one class dominates---both failure modes we document in
\S\ref{sec:results-detect}. Items are the independent unit throughout, and every confidence
interval resamples items rather than verification calls.

% \section{Submission of papers to NeurIPS 2026}


% Please read the instructions below carefully and follow them faithfully.


% % \subsection{Style}


% % Papers to be submitted to NeurIPS 2026 must be prepared according to the
% % instructions presented here. Papers may only be up to {\bf nine} pages long, including figures. \textbf{Papers that exceed the page limit will not be reviewed (or in any other way considered) for presentation at the conference.}
% % Additional pages \emph{containing acknowledgments, references, checklist, and optional technical appendices} do not count as content pages. 


% % The margins in 2026 are the same as those in previous years.


% % Authors are required to use the NeurIPS \LaTeX{} style files obtainable at the NeurIPS website as indicated below. Please make sure you use the current files and not previous versions. Tweaking the style files may be grounds for desk rejection.


% % \subsection{Retrieval of style files}


% % The style files for NeurIPS and other conference information are available on the website at
% % \begin{center}
% %   \url{https://neurips.cc}.
% % \end{center}
% % % The file \verb+neurips_2026.pdf+ contains these instructions and illustrates the various formatting requirements your NeurIPS paper must satisfy.


% % The only supported style file for NeurIPS 2026 is \verb+neurips_2026.sty+, rewritten for \LaTeXe{}. \textbf{Previous style files for \LaTeX{} 2.09,  Microsoft Word, and RTF are no longer supported.}


% % The \LaTeX{} style file contains three optional arguments: 
% % \begin{itemize}
% % \item \verb+final+, which creates a camera-ready copy, 
% % \item \verb+preprint+, which creates a preprint for submission to, e.g., arXiv, \item \verb+nonatbib+, which will not load the \verb+natbib+ package for you in case of package clash.
% % \end{itemize}


% % \paragraph{Preprint option}
% % If you wish to post a preprint of your work online, e.g., on arXiv, using the NeurIPS style, please use the \verb+preprint+ option. This will create a nonanonymized version of your work with the text ``Preprint. Work in progress.'' in the footer. This version may be distributed as you see fit, as long as you do not say which conference it was submitted to. Please \textbf{do not} use the \verb+final+ option, which should \textbf{only} be used for papers accepted to NeurIPS.


% % At submission time, please omit the \verb+final+ and \verb+preprint+ options. This will anonymize your submission and add line numbers to aid review. Please do \emph{not} refer to these line numbers in your paper as they will be removed during generation of camera-ready copies.


% % The file \verb+neurips_2026.tex+ may be used as a ``shell'' for writing your paper. All you have to do is replace the author, title, abstract, and text of the paper with your own.


% % The formatting instructions contained in these style files are summarized in Sections \ref{gen_inst}, \ref{headings}, and \ref{others} below.


% % \section{General formatting instructions}
% % \label{gen_inst}


% % The text must be confined within a rectangle 5.5~inches (33~picas) wide and
% % 9~inches (54~picas) long. The left margin is 1.5~inch (9~picas).  Use 10~point
% % type with a vertical spacing (leading) of 11~points.  Times New Roman is the
% % preferred typeface throughout, and will be selected for you by default.
% % Paragraphs are separated by \nicefrac{1}{2}~line space (5.5 points), with no
% % indentation.


% % The paper title should be 17~point, initial caps/lower case, bold, centered
% % between two horizontal rules. The top rule should be 4~points thick and the
% % bottom rule should be 1~point thick. Allow \nicefrac{1}{4}~inch space above and
% % below the title to rules. All pages should start at 1~inch (6~picas) from the
% % top of the page.


% % For the final version, authors' names are set in boldface, and each name is
% % centered above the corresponding address. The lead author's name is to be listed
% % first (left-most), and the co-authors' names (if different address) are set to
% % follow. If there is only one co-author, list both author and co-author side by
% % side.


% % Please pay special attention to the instructions in Section \ref{others}
% % regarding figures, tables, acknowledgments, and references.

% % \section{Headings: first level}
% % \label{headings}


% % All headings should be lower case (except for first word and proper nouns),
% % flush left, and bold.


% % First-level headings should be in 12-point type.


% % \subsection{Headings: second level}


% % Second-level headings should be in 10-point type.


% % \subsubsection{Headings: third level}


% % Third-level headings should be in 10-point type.


% % \paragraph{Paragraphs}


% % There is also a \verb+\paragraph+ command available, which sets the heading in
% % bold, flush left, and inline with the text, with the heading followed by 1\,em
% % of space.


% % \section{Citations, figures, tables, references}
% % \label{others}


% % These instructions apply to everyone.


% % \subsection{Citations within the text}


% % The \verb+natbib+ package will be loaded for you by default.  Citations may be
% % author/year or numeric, as long as you maintain internal consistency.  As to the
% % format of the references themselves, any style is acceptable as long as it is
% % used consistently.


% % The documentation for \verb+natbib+ may be found at
% % \begin{center}
% %   \url{http://mirrors.ctan.org/macros/latex/contrib/natbib/natnotes.pdf}
% % \end{center}
% % Of note is the command \verb+\citet+, which produces citations appropriate for
% % use in inline text.  For example,
% % \begin{verbatim}
% %    \citet{hasselmo} investigated\dots
% % \end{verbatim}
% % produces
% % \begin{quote}
% %   Hasselmo, et al.\ (1995) investigated\dots
% % \end{quote}


% % If you wish to load the \verb+natbib+ package with options, you may add the
% % following before loading the \verb+neurips_2026+ package:
% % \begin{verbatim}
% %    \PassOptionsToPackage{options}{natbib}
% % \end{verbatim}


% % If \verb+natbib+ clashes with another package you load, you can add the optional
% % argument \verb+nonatbib+ when loading the style file:
% % \begin{verbatim}
% %    \usepackage[nonatbib]{neurips_2026}
% % \end{verbatim}


% % As submission is double blind, refer to your own published work in the third
% % person. That is, use ``In the previous work of Jones et al.\ [4],'' not ``In our
% % previous work [4].'' If you cite your other papers that are not widely available
% % (e.g., a journal paper under review), use anonymous author names in the
% % citation, e.g., an author of the form ``A.\ Anonymous'' and include a copy of the anonymized paper in the supplementary material.


% % \subsection{Footnotes}


% % Footnotes should be used sparingly.  If you do require a footnote, indicate
% % footnotes with a number\footnote{Sample of the first footnote.} in the
% % text. Place the footnotes at the bottom of the page on which they appear.
% % Precede the footnote with a horizontal rule of 2~inches (12~picas).


% % Note that footnotes are properly typeset \emph{after} punctuation
% % marks.\footnote{As in this example.}


% % \subsection{Figures}


% % \begin{figure}
% %   \centering
% %   \fbox{\rule[-.5cm]{0cm}{4cm} \rule[-.5cm]{4cm}{0cm}}
% %   \caption{Sample figure caption. Explain what the figure shows and add a key take-away message to the caption.}
% % \end{figure}


% % All artwork must be neat, clean, and legible. Lines should be dark enough for
% %  reproduction purposes. The figure number and caption always appear after the
% % figure. Place one line space before the figure caption and one line space after
% % the figure. The figure caption should be lower case (except for the first word and proper nouns); figures are numbered consecutively.


% % You may use color figures.  However, it is best for the figure captions and the
% % paper body to be legible if the paper is printed in either black/white or in
% % color.


% % \subsection{Tables}


% % All tables must be centered, neat, clean, and legible.  The table number and
% % title always appear before the table.  See Table~\ref{sample-table}.


% % Place one line space before the table title, one line space after the
% % table title, and one line space after the table. The table title must
% % be lower case (except for the first word and proper nouns); tables are
% % numbered consecutively.


% % Note that publication-quality tables \emph{do not contain vertical rules}. We
% % strongly suggest the use of the \verb+booktabs+ package, which allows for
% % typesetting high-quality, professional tables:
% % \begin{center}
% %   \url{https://www.ctan.org/pkg/booktabs}
% % \end{center}
% % This package was used to typeset Table~\ref{sample-table}.


% % \begin{table}
% %   \caption{Sample table caption. Explain what the table shows and add a key take-away message to the caption.}
% %   \label{sample-table}
% %   \centering
% %   \begin{tabular}{lll}
% %     \toprule
% %     \multicolumn{2}{c}{Part}                   \\
% %     \cmidrule(r){1-2}
% %     Name     & Description     & Size ($\mu$m) \\
% %     \midrule
% %     Dendrite & Input terminal  & $\approx$100     \\
% %     Axon     & Output terminal & $\approx$10      \\
% %     Soma     & Cell body       & up to $10^6$  \\
% %     \bottomrule
% %   \end{tabular}
% % \end{table}

% % \subsection{Math}
% % Note that display math in bare TeX commands will not create correct line numbers for submission. Please use LaTeX (or AMSTeX) commands for unnumbered display math. (You really shouldn't be using \$\$ anyway; see \url{https://tex.stackexchange.com/questions/503/why-is-preferable-to} and \url{https://tex.stackexchange.com/questions/40492/what-are-the-differences-between-align-equation-and-displaymath} for more information.)

% % \subsection{Final instructions}

% % Do not change any aspects of the formatting parameters in the style files.  In
% % particular, do not modify the width or length of the rectangle the text should
% % fit into, and do not change font sizes. Please note that pages should be
% % numbered.


% % \section{Preparing PDF files}


% % Please prepare submission files with paper size ``US Letter,'' and not, for
% % example, ``A4.''


% % Fonts were the main cause of problems in the past years. Your PDF file must only
% % contain Type 1 or Embedded TrueType fonts. Here are a few instructions to
% % achieve this.


% % \begin{itemize}


% % \item You should directly generate PDF files using \verb+pdflatex+.


% % \item You can check which fonts a PDF files uses.  In Acrobat Reader, select the
% %   menu Files$>$Document Properties$>$Fonts and select Show All Fonts. You can
% %   also use the program \verb+pdffonts+ which comes with \verb+xpdf+ and is
% %   available out-of-the-box on most Linux machines.


% % \item \verb+xfig+ ``patterned'' shapes are implemented with bitmap fonts.  Use
% %   "solid" shapes instead.


% % \item The \verb+\bbold+ package almost always uses bitmap fonts.  You should use
% %   the equivalent AMS Fonts:
% % \begin{verbatim}
% %    \usepackage{amsfonts}
% % \end{verbatim}
% % followed by, e.g., \verb+\mathbb{R}+, \verb+\mathbb{N}+, or \verb+\mathbb{C}+
% % for $\mathbb{R}$, $\mathbb{N}$ or $\mathbb{C}$.  You can also use the following
% % workaround for reals, natural and complex:
% % \begin{verbatim}
% %    \newcommand{\RR}{I\!\!R} %real numbers
% %    \newcommand{\Nat}{I\!\!N} %natural numbers
% %    \newcommand{\CC}{I\!\!\!\!C} %complex numbers
% % \end{verbatim}
% % Note that \verb+amsfonts+ is automatically loaded by the \verb+amssymb+ package.


% % \end{itemize}


% % If your file contains type 3 fonts or non embedded TrueType fonts, we will ask
% % you to fix it.


% % \subsection{Margins in \LaTeX{}}


% % Most of the margin problems come from figures positioned by hand using
% % \verb+\special+ or other commands. We suggest using the command
% % \verb+\includegraphics+ from the \verb+graphicx+ package. Always specify the
% % figure width as a multiple of the line width as in the example below:
% % \begin{verbatim}
% %    \usepackage[pdftex]{graphicx} ...
% %    \includegraphics[width=0.8\linewidth]{myfile.pdf}
% % \end{verbatim}
% % See Section 4.4 in the graphics bundle documentation
% % (\url{http://mirrors.ctan.org/macros/latex/required/graphics/grfguide.pdf})


% % A number of width problems arise when \LaTeX{} cannot properly hyphenate a
% % line. Please give LaTeX hyphenation hints using the \verb+\-+ command when
% % necessary.

% % \begin{ack}
% % Use unnumbered first level headings for the acknowledgments. All acknowledgments
% % go at the end of the paper before the list of references. Moreover, you are required to declare
% % funding (financial activities supporting the submitted work) and competing interests (related financial activities outside the submitted work).
% % More information about this disclosure can be found at: \url{https://neurips.cc/Conferences/2026/PaperInformation/FundingDisclosure}.


% % Do {\bf not} include this section in the anonymized submission, only in the final paper. You can use the \texttt{ack} environment provided in the style file to automatically hide this section in the anonymized submission.
% % \end{ack}

\section*{References}


References follow the acknowledgments in the camera-ready paper. Use unnumbered first-level heading for
the references. Any choice of citation style is acceptable as long as you are
consistent. It is permissible to reduce the font size to \verb+small+ (9 point)
when listing the references.
Note that the Reference section does not count towards the page limit.
\medskip


{
\small


[1] Alexander, J.A.\ \& Mozer, M.C.\ (1995) Template-based algorithms for
connectionist rule extraction. In G.\ Tesauro, D.S.\ Touretzky and T.K.\ Leen
(eds.), {\it Advances in Neural Information Processing Systems 7},
pp.\ 609--616. Cambridge, MA: MIT Press.


[2] Bower, J.M.\ \& Beeman, D.\ (1995) {\it The Book of GENESIS: Exploring
  Realistic Neural Models with the GEneral NEural SImulation System.}  New York:
TELOS/Springer--Verlag.


[3] Hasselmo, M.E., Schnell, E.\ \& Barkai, E.\ (1995) Dynamics of learning and
recall at excitatory recurrent synapses and cholinergic modulation in rat
hippocampal region CA3. {\it Journal of Neuroscience} {\bf 15}(7):5249-5262.
}


%%%%%%%%%%%%%%%%%%%%%%%%%%%%%%%%%%%%%%%%%%%%%%%%%%%%%%%%%%%%

\appendix

\section{Technical appendices and supplementary material}
Technical appendices with additional results, figures, graphs, and proofs may be submitted with the paper submission before the full submission deadline (see above). You can upload a ZIP file for videos or code, but do not upload a separate PDF file for the appendix. There is no page limit for the technical appendices. 

Note: Think of the appendix as ``optional reading'' for reviewers. The paper must be able to stand alone without the appendix; for example, adding critical experiments that support the main claims to an appendix is inappropriate. 

%%%%%%%%%%%%%%%%%%%%%%%%%%%%%%%%%%%%%%%%%%%%%%%%%%%%%%%%%%%%

\newpage
\input{checklist.tex}


\end{document}
